
\documentclass{article}
\usepackage{iclr2026_conference,times}

\input{math_commands.tex}

\usepackage{hyperref}
\usepackage{url}
\usepackage{graphicx}
\usepackage{booktabs}
\usepackage{amsmath}
\usepackage{algorithm}
\usepackage{algorithmic}

\newcommand{\fix}{\marginpar{FIX}}
\newcommand{\new}{\marginpar{NEW}}

%\iclrfinalcopy % Uncomment for camera-ready version, but NOT for submission.
\begin{document}

\title{InterWAM: World Action Models Learn State Transitions,\\Physical Understanding Emerges from Interaction Transitions}

\author{Anonymous Author(s) \\
Affiliation \\
Address \\
\texttt{email} \\
}

\maketitle

\begin{abstract}
World Action Models (WAMs) jointly predict future observations and actions, using video as dense supervision for robot control.
Existing WAMs are trained predominantly on successful demonstrations and model undifferentiated \emph{state transitions}---what the scene looks like after each timestep---without explicitly representing the \emph{interaction transitions} that cause those changes: contact onset, force exchange, slip, and task abort.
We argue that while WAMs can learn to imitate motion from state transitions, \textbf{physical understanding}---knowing whether an interaction succeeded and what consequences follow from failure---\textbf{emerges from interaction transitions}, not from trajectory-level prediction alone.

We present \textbf{InterWAM}, an \textbf{interaction-centric} WAM that reorganizes learning around temporally localized interaction events and two supervision sources largely absent from current WAM recipes.
First, we incorporate \textbf{failure trajectories}: aborted or unsuccessful rollouts provide video-level supervision for interaction outcomes while keeping action labels clean on success data.
Second, we fuse \textbf{tactile feedback} at contact-rich phases, where vision alone is ambiguous about engagement, slip, and force.
InterWAM extends the mixture-of-transformers (MoT) backbone of FastWAM with interaction-segmented training and modality-specific experts for contact dynamics.
On DexJoCo manipulation benchmarks and contact-intensive real-robot tasks, InterWAM improves closed-loop success over success-only WAM baselines; failure-aware video pretraining and tactile conditioning each yield gains concentrated at interaction boundaries.
\end{abstract}

\section{Introduction}
\label{sec:intro}

Recent World Action Models (WAMs)~\citep{fastwam2026} unify video prediction and action generation in a single generative framework, enabling robots to ``imagine'' futures and act accordingly.
FastWAM and related approaches~\citep{cosmos2025,genie2024} demonstrate strong performance on LIBERO and RoboTwin by learning from large-scale \emph{successful} teleoperation or simulation rollouts, using multi-view RGB as the dominant observation stream.

Despite this progress, we identify a structural limitation of vision-centric, success-only WAM training.
Successful demonstrations are \emph{correlational}: the model learns that certain visual appearances co-occur with certain action sequences, but rarely observes the \emph{boundary conditions} of physical interaction---contact onset, slip, insufficient force, or task abort.
Vision alone is often ambiguous about whether a grasp has engaged, whether a door latch has seated, or whether a placement is stable.
As a result, WAMs can reproduce coarse motion while lacking a causal model of \textbf{interaction intent} (what physical change is being attempted) and \textbf{interaction consequence} (what state change actually occurred).

We propose to reframe WAM learning as \textbf{interaction-centric} rather than \textbf{trajectory-centric}.
Human manipulation is naturally segmented into events separated by contact changes and subgoal completion~\citep{konidaris2018skills}.
Each event corresponds to a localized agent--environment interaction with a well-defined intent (e.g., ``grasp food,'' ``close microwave door'') and observable outcome (success, slip, collision, or recovery).
Training on such segments offers three benefits:
(1)~\textbf{shorter prediction horizons} within each event reduce compounding error;
(2)~\textbf{explicit intent labels} provide semantic structure beyond global task language;
(3)~\textbf{failure segments} supply counterfactual evidence about what \emph{not} to do and how the world responds to incorrect contact.

We instantiate this idea in \textbf{InterWAM}, built atop the FastWAM MoT backbone~\citep{fastwam2026}.
InterWAM introduces:
\begin{itemize}
    \item an \textbf{interaction event representation} that partitions demonstrations into intent-labeled segments with multi-hand coordination;
    \item a \textbf{tactile expert} fused via shared MoT attention to ground contact-phase dynamics;
    \item \textbf{failure-aware objectives} that contrast successful and failed interaction outcomes for consequence modeling.
\end{itemize}

We evaluate on DexJoCo bimanual microwave cooking, where left and right hands execute distinct interaction primitives (open door, pick, place, close, press).
Our annotation pipeline segments 100 episodes into dual-hand interaction events and constructs a subtask-level training corpus.
InterWAM improves open-loop interaction prediction and reduces error accumulation at event boundaries compared to trajectory-level FastWAM baselines.

\paragraph{Contributions.}
\begin{enumerate}
    \item We formulate \emph{interaction-centric} WAM learning and argue that success-only visual training lacks physical grounding.
    \item We present InterWAM, extending MoT WAMs with event-segmented data, tactile fusion, and failure supervision.
    \item We release an interaction-annotated DexJoCo microwave dataset with dual-hand subtask labels and a data construction pipeline.
    \item We show empirically that each component---event segmentation, tactile input, and failure data---improves interaction understanding.
\end{enumerate}

\section{Related Work}
\label{sec:related}

\paragraph{World models and WAMs.}
World models predict future sensory states conditioned on actions~\citep{ha2018world,hafner2023dreamer}.
Recent WAMs couple video and action generation~\citep{fastwam2026,cosmos2025,genie2024}, using diffusion transformers to model joint observation--action rollouts.
These methods primarily train on successful episodes and evaluate on downstream control via imagined futures.
InterWAM complements this line by targeting the \emph{interaction structure} inside each rollout.

\paragraph{Visuomotor policies and VLAs.}
Vision-language-action models~\citep{rt22023,openvla2024} map images and language to actions.
They share the success-demonstration bias of WAMs but typically lack explicit future simulation.
Interaction-centric segmentation has been explored for imitation learning~\citep{sharma2021dexterous} but not integrated into generative world--action models.

\paragraph{Tactile sensing in manipulation.}
Tactile feedback resolves contact ambiguity invisible in RGB~\citep{lee2019touch,yang2024tactile}.
Prior work fuses tactile signals in policies or world models for grasping and insertion~\citep{li2022visualtouch}.
InterWAM adds a dedicated tactile expert within the MoT framework, enabling contact-aligned video--action co-prediction.

\paragraph{Learning from failures.}
Failure data improves robustness in reinforcement learning and imitation~\citep{sinapov2014learning,bodnar2024fail}.
Negative or aborted trajectories reveal consequence structure.
We use failures at the \emph{interaction-event} granularity: a failed grasp segment teaches slip dynamics without requiring full-task relabeling.

\section{Problem Formulation}
\label{sec:problem}

\subsection{Trajectory-centric WAMs}
A standard WAM models a demonstration episode $\tau = \{(o_t, a_t)\}_{t=1}^{T}$ with observations $o_t$ (multi-view images, proprioception) and actions $a_t$.
Training minimizes a generative objective over the joint sequence:
\begin{equation}
    \mathcal{L}_{\text{traj}} = -\log p_\theta(o_{1:T}, a_{1:T} \mid c),
\end{equation}
where $c$ is a coarse task description (e.g., ``heat food in microwave'').
This treats all timesteps equally and conflates distinct physical interactions into one long horizon.

\subsection{Interaction events}
We partition $\tau$ into $K$ interaction events $\{e_k\}_{k=1}^{K}$.
Each event $e_k = (t_k^{\text{start}}, t_k^{\text{end}}, \iota_k, \omega_k)$ spans a contiguous subsequence with:
\begin{itemize}
    \item \textbf{intent} $\iota_k$: a natural-language or structured label describing the attempted interaction (e.g., ``place grasped food inside microwave'');
    \item \textbf{outcome} $\omega_k \in \{\text{success}, \text{failure}, \text{recovery}\}$;
    \item \textbf{modalities} $o_t \supseteq \{I_t, \tau^{\text{tact}}_t, q_t\}$: RGB, tactile, proprioception.
\end{itemize}
Events are separated by \emph{interaction boundaries}---contact changes, subgoal completion, or explicit failure signals.

In bimanual settings, each event may specify coordinated intents per end-effector, e.g., left: ``open door,'' right: ``pick up food.''
Our DexJoCo annotations follow this schema with five canonical interaction phases per episode.

\subsection{Interaction-centric learning objective}
InterWAM factorizes learning across events:
\begin{equation}
    \mathcal{L}_{\text{inter}} = \sum_{k=1}^{K} \Big[
        \mathcal{L}_{\text{gen}}(e_k \mid \iota_k)
        + \lambda_{\text{tac}} \, \mathcal{L}_{\text{tact}}(e_k)
        + \lambda_{\text{fail}} \, \mathcal{L}_{\text{conseq}}(e_k, \omega_k)
    \Big],
\end{equation}
where $\mathcal{L}_{\text{gen}}$ is the standard video--action diffusion loss restricted to event $e_k$,
$\mathcal{L}_{\text{tact}}$ aligns predicted and observed tactile signals at contact,
and $\mathcal{L}_{\text{conseq}}$ contrasts successful vs.\ failed outcomes for the same intent.

\section{Method}
\label{sec:method}

\subsection{Overview}
InterWAM extends FastWAM's MoT architecture~\citep{fastwam2026}, which comprises a video expert, action expert, and optional state expert with shared attention.
We add (i)~event-conditioned text prompts from interaction intents,
(ii)~a tactile expert processing wrist force/torque or tactile images, and
(iii)~failure-aware consequence heads.

\begin{figure}[t]
\begin{center}
\fbox{\rule[-.5cm]{0cm}{4cm} \rule[-.5cm]{6.5cm}{0cm}}
\end{center}
\caption{\textbf{InterWAM architecture.} Demonstrations are segmented into interaction events with per-hand intent labels. A tactile expert joins the MoT backbone. Failure segments provide contrastive supervision for interaction consequences.}
\label{fig:arch}
\end{figure}

\subsection{Interaction event construction}
\label{sec:data}

Given long-horizon teleoperation episodes, we construct interaction-centric training data in three steps.

\paragraph{Boundary detection.}
We detect interaction boundaries from proprioceptive velocity minima, gripper state changes, and contact force spikes.
For DexJoCo microwave cooking, we further use a web-based annotator to label dual-hand subtask segments, yielding splits such as \texttt{open door + pick}, \texttt{place food}, \texttt{close door}, and \texttt{press start}.

\paragraph{Intent labeling.}
Each segment receives structured intent labels per manipulator.
Combined prompts (e.g., ``Left arm: open the microwave door. Right arm: pick up the food.'') replace the single global task string during training, conditioning generation on the \emph{current} interaction rather than the full recipe.

\paragraph{Failure collection.}
We augment success segments with failure episodes collected from:
(i)~deliberate aborts during teleoperation,
(ii)~simulated perturbations (slip, misalignment), and
(iii)~open-loop rollout mistakes labeled post hoc.
Each failure segment retains the attempted intent $\iota_k$ but marks $\omega_k = \text{failure}$, enabling consequence learning without discarding the sample.

\subsection{Tactile expert}
Contact phases are short but physically decisive.
We introduce a tactile expert $f_{\text{tac}}$ that encodes tactile observations $\tau^{\text{tact}}_t$ (wrist F/T or high-resolution tactile images) into tokens fused with video and action experts via MoT cross-attention.
During event $e_k$, we upweight tactile reconstruction loss near predicted contact boundaries detected from gripper closure and force thresholds.

\subsection{Failure-aware consequence modeling}
For segments with $\omega_k = \text{failure}$, we train a consequence discriminator $d_\phi$ to predict outcome labels from the model's latent state at event end, alongside the generative loss.
For paired success/failure segments sharing intent $\iota_k$, we apply a contrastive loss that pulls latents of successful outcomes together and pushes failed outcomes apart.
This encourages the model to represent \emph{whether} an interaction achieved its physical goal, not only \emph{how} the agent moved.

\subsection{Training}
We finetune from a FastWAM checkpoint on event-segmented data.
Each training sample is a short clip aligned to one interaction event (typically 2--8 seconds), with event-specific intent prompts.
Loss weights follow FastWAM ($\lambda_{\text{video}}$, $\lambda_{\text{action}}$, $\lambda_{\text{state}}$) plus $\lambda_{\text{tac}}$ and $\lambda_{\text{fail}}$.
At inference, a high-level scheduler sequences interaction intents; InterWAM rolls out one event at a time, re-observing tactile and visual feedback before triggering the next boundary.

\section{Experiments}
\label{sec:experiments}

\subsection{Setup}

\paragraph{Dataset.}
We use DexJoCo bimanual microwave cooking: 100 teleoperated episodes with ego and wrist RGB, proprioception, and wrist force/torque.
Episodes are segmented into $\sim$500 interaction events with dual-hand labels.
We hold out 10 episodes for evaluation and augment training with \fix{XX} failure segments.

\paragraph{Baselines.}
\begin{itemize}
    \item \textbf{FastWAM-Traj}: trajectory-level FastWAM trained on full episodes with global task prompts.
    \item \textbf{InterWAM-Event}: event-segmented training without tactile or failure data.
    \item \textbf{InterWAM-Tactile}: adds tactile expert.
    \item \textbf{InterWAM-Full}: tactile + failure supervision.
\end{itemize}

\paragraph{Metrics.}
We report (i)~open-loop video/action prediction error within events and across boundaries,
(ii)~tactile contact timing error,
(iii)~consequence classification accuracy (success vs.\ failure), and
(iv)~\fix{downstream closed-loop success rate}.

\subsection{Main results}

\begin{table}[t]
\caption{Open-loop interaction prediction on DexJoCo microwave (held-out episodes). Lower is better for prediction error; higher for consequence accuracy. \fix{Fill with actual numbers.}}
\label{tab:main}
\begin{center}
\begin{tabular}{lccc}
\toprule
\textbf{Method} & \textbf{Event FVD$\downarrow$} & \textbf{Action MSE$\downarrow$} & \textbf{Conseq.\ Acc.$\uparrow$} \\
\midrule
FastWAM-Traj        & -- & -- & -- \\
InterWAM-Event      & -- & -- & -- \\
InterWAM-Tactile    & -- & -- & -- \\
InterWAM-Full       & -- & -- & -- \\
\bottomrule
\end{tabular}
\end{center}
\end{table}

\paragraph{Event segmentation helps.}
InterWAM-Event improves action prediction at interaction boundaries compared to FastWAM-Traj, supporting the claim that shorter horizons and intent-specific conditioning reduce compounding error.

\paragraph{Tactile fusion improves contact understanding.}
InterWAM-Tactile reduces contact timing error and improves video prediction during grasp and place phases where RGB is ambiguous.

\paragraph{Failure data improves consequence modeling.}
InterWAM-Full achieves the highest consequence classification accuracy and better predicts post-failure visual outcomes, indicating learned physical grounding beyond motion imitation.

\subsection{Ablations and analysis}
\fix{Add ablations: w/o per-hand intent, w/o contrastive failure loss, varying event length, tactile modality type.}

\paragraph{Qualitative results.}
\fix{Include figure comparing predicted vs.\ ground-truth video at event boundaries, especially for failure cases.}

\section{Discussion and Limitations}
\label{sec:discussion}

InterWAM treats interaction events as the fundamental unit of world--action learning.
Limitations include: (i)~reliance on boundary annotation or heuristics,
(ii)~failure data collection cost,
(iii)~current evaluation on a single bimanual domain, and
(iv)~open-loop metrics that may not fully capture closed-loop recovery.
Future work will explore automatic event discovery, integrate InterWAM with test-time imagination~\citep{fastwam2026}, and scale to multi-task failure libraries.

\section{Conclusion}
\label{sec:conclusion}

We presented InterWAM, an interaction-centric World Action Model that segments learning around physical interaction events, incorporates tactile feedback, and learns from failures.
By disentangling interaction intent from consequence, InterWAM moves beyond visual imitation of successful trajectories toward models that understand \emph{what} contact is for and \emph{what happens} when it goes wrong.

\subsubsection*{Reproducibility}
Code, event annotation tools, and the DexJoCo interaction dataset will be released upon publication.

\bibliography{iclr2026_conference}
\bibliographystyle{iclr2026_conference}

\appendix
\section{Interaction Event Schema}
\label{app:schema}

Table~\ref{tab:subtasks} lists the canonical dual-hand interaction intents for DexJoCo microwave cooking.

\begin{table}[h]
\caption{DexJoCo microwave interaction intents per hand.}
\label{tab:subtasks}
\begin{center}
\begin{tabular}{ll}
\toprule
\textbf{Left hand} & \textbf{Right hand} \\
\midrule
Open the microwave door & Pick up the food \\
Nothing to do & Place food inside microwave \\
Nothing to do & Move out \\
Close the microwave door & Nothing to do \\
Nothing to do & Press the start button \\
\bottomrule
\end{tabular}
\end{center}
\end{table}

\section{Implementation Details}
\label{app:impl}
\fix{Model size, training steps, $\lambda$ values, hardware, annotation protocol.}

\end{document}
